% ============================================================================
%  Experimental Results + Discussion (DROP-IN REPLACEMENT)
%
%  HOW TO USE
%  ----------
%  In main.tex, DELETE the current Section "Experimental Results"
%  (the block from "\section{Experimental Results}" down to just before
%  "\section{Conclusion}") and replace it with one line:
%
%      % ============================================================================
%  Experimental Results + Discussion (DROP-IN REPLACEMENT)
%
%  HOW TO USE
%  ----------
%  In main.tex, DELETE the current Section "Experimental Results"
%  (the block from "\section{Experimental Results}" down to just before
%  "\section{Conclusion}") and replace it with one line:
%
%      % ============================================================================
%  Experimental Results + Discussion (DROP-IN REPLACEMENT)
%
%  HOW TO USE
%  ----------
%  In main.tex, DELETE the current Section "Experimental Results"
%  (the block from "\section{Experimental Results}" down to just before
%  "\section{Conclusion}") and replace it with one line:
%
%      % ============================================================================
%  Experimental Results + Discussion (DROP-IN REPLACEMENT)
%
%  HOW TO USE
%  ----------
%  In main.tex, DELETE the current Section "Experimental Results"
%  (the block from "\section{Experimental Results}" down to just before
%  "\section{Conclusion}") and replace it with one line:
%
%      \input{results_section}
%
%  Requires (already in main.tex preamble): booktabs, amssymb,
%  \xhdr command, cleveref. Uses table* (full page width) for the wide
%  statistics table, so it must be in a twocolumn document (it is).
%
%  All numbers below are produced by:
%      python eval/statistical_tests.py   -> eval/statistical_tests.json
%      python eval/failure_analysis.py    -> eval/failure_analysis.json
%      python eval/compare_scores.py      -> eval/comparison_summary.json
%  in the ai-interview-coach repository. Regenerate those if data changes.
% ============================================================================

\section{Experimental Results}
\label{sec:results}

\xhdr{Evaluation protocol.} We evaluate both models on a fixed held-out test
set of 200 examples. The test questions and the simulated student answers are
identical for the baseline and the trained model, so any score difference is
attributable only to the coaching model. For each example, the judge model
(Qwen3.5-9B) scores the feedback on five dimensions (Technical Correctness,
Specificity, Helpfulness, Actionability, Interview Coaching Quality), each on a
1--20 scale; the overall score is their mean. We report per-dimension means and,
because the same 200 items are scored by both models, we treat the two score
vectors as \emph{paired} samples and test the difference with a paired
$t$-test, the Wilcoxon signed-rank test (robust to the non-normal,
ceiling-bound score distribution), and an exact binomial sign test on the
non-tie items. We also report Cohen's $d$ for paired samples as an effect size.

\xhdr{Aggregate results.} As shown in \cref{tab:results}, the SFT+DPO model
attains a marginally higher mean overall score than the baseline
($17.03$ vs.\ $16.86$, $\Delta=+0.17$ on the 1--20 scale). The improvement is
consistent in direction across all five dimensions, but small in magnitude
(per-dimension $\Delta$ between $+0.14$ and $+0.19$).

\xhdr{Statistical significance.} The improvement is \emph{not} statistically
significant. For the overall score, the paired $t$-test gives $t=0.81$,
$p=0.42$, and the Wilcoxon signed-rank test gives $p=0.70$; the effect size is
negligible (Cohen's $d=0.057$). No individual dimension reaches significance
(all $p>0.38$; see \cref{tab:stats}). On an item-by-item basis the trained
model wins on 48 examples, loses on 40, and ties on 112; the sign test over the
88 non-tie items is also non-significant ($p=0.46$). We therefore cannot reject
the null hypothesis that SFT+DPO leaves feedback quality unchanged under this
judge and rubric.

\begin{table}[t]
  \centering
  \small
  \begin{tabular}{@{}lccc@{}}
    \toprule
    Dimension & Baseline & SFT+DPO & $\Delta$ \\
    \midrule
    Technical Correctness        & 16.79 & 16.92 & $+0.14$ \\
    Specificity                  & 16.88 & 17.06 & $+0.19$ \\
    Helpfulness                  & 16.88 & 17.05 & $+0.18$ \\
    Actionability                & 16.89 & 17.05 & $+0.17$ \\
    Interview Coaching Quality   & 16.88 & 17.06 & $+0.19$ \\
    \midrule
    \textbf{Overall Score}       & \textbf{16.86} & \textbf{17.03} & $\mathbf{+0.17}$ \\
    \bottomrule
  \end{tabular}
  \caption{Mean judge scores (1--20) on the 200-example test set. The trained
  model is consistently but marginally higher across all dimensions.}
  \label{tab:results}
\end{table}

\begin{table*}[t]
  \centering
  \small
  \begin{tabular}{@{}lcccccc@{}}
    \toprule
    Dimension & Baseline (SD) & SFT+DPO (SD) & $\Delta$ & Paired $t$ ($p$) & Wilcoxon $p$ & Cohen's $d$ \\
    \midrule
    Technical Correctness        & 16.79 (2.40) & 16.92 (1.92) & $+0.14$ & $0.63\ (0.53)$ & $0.73$ & $0.044$ \\
    Specificity                  & 16.88 (2.37) & 17.06 (1.82) & $+0.19$ & $0.88\ (0.38)$ & $0.65$ & $0.062$ \\
    Helpfulness                  & 16.88 (2.37) & 17.05 (1.85) & $+0.18$ & $0.83\ (0.41)$ & $0.69$ & $0.058$ \\
    Actionability                & 16.89 (2.38) & 17.05 (1.84) & $+0.17$ & $0.78\ (0.44)$ & $0.81$ & $0.055$ \\
    Interview Coaching Quality   & 16.88 (2.37) & 17.06 (1.82) & $+0.19$ & $0.88\ (0.38)$ & $0.69$ & $0.062$ \\
    \midrule
    \textbf{Overall Score}       & \textbf{16.86 (2.37)} & \textbf{17.03 (1.85)} & $\mathbf{+0.17}$ & $\mathbf{0.81\ (0.42)}$ & $\mathbf{0.70}$ & $\mathbf{0.057}$ \\
    \bottomrule
  \end{tabular}
  \caption{Paired statistical comparison of baseline vs.\ SFT+DPO over the 200
  test examples. No dimension reaches significance at $\alpha=0.05$, and all
  effect sizes are negligible ($d<0.07$). Win/loss/tie on the overall score is
  48/40/112; the sign test over non-tie items gives $p=0.46$.}
  \label{tab:stats}
\end{table*}

\section{Discussion}
\label{sec:discussion}

The absence of a significant difference is itself an informative result. We
identify three contributing factors and support them with a per-item failure
analysis.

\xhdr{Ceiling effect.} The judge already rates the baseline very highly,
leaving little headroom for improvement. 65\% of baseline outputs score
$\geq 18/20$, and on 191 of 200 items \emph{both} models score $\geq 15$. When
the starting point is already near the top of the scale, even a genuinely
better model cannot move the mean by much.

\xhdr{Coarse judge resolution.} The judge's scores are effectively discrete,
clustering on a handful of values (predominantly 15 and 18), and the five
rubric dimensions almost always move together (their per-item scores are nearly
identical). This low resolution injects noise and masks the fine-grained
differences that distinguish good coaching feedback from excellent coaching
feedback.

\xhdr{Offsetting gains and regressions.} The flat average hides real movement
in both directions: 48 items improve, 40 regress, and these roughly cancel. The
trained model's largest wins are recoveries of \emph{catastrophic} baseline
failures---three items jump from 1 to 18 (\texttt{test\_0001},
\texttt{test\_0073}, \texttt{test\_0119}) by removing a baseline hallucination
that claimed the student had missed a concept actually present in the answer.
However, the trained model also introduces \emph{new} technical errors that the
baseline did not make: in the most severe regression, \texttt{test\_0194} drops
from 18 to 1 because the feedback falsely asserts that Java's
\texttt{@Override} annotation throws a runtime \texttt{java.lang.Error}.
Categorizing the judge's stated reason on the 40 regressions, 4 are clear
technical inaccuracies and the remainder are smaller, mostly within the judge's
own scoring noise (e.g.\ slightly less complete or less specific feedback).

\xhdr{Test-set diversity limitation.} A confound limits what this experiment can
demonstrate. Although our methodology defines five student-answer types
(correct-but-incomplete, partially correct, incorrect, too vague, verbose but
unfocused), in the released data all 200 test items (and all training items)
carry the label \texttt{partially\_correct}, and the question
category/difficulty metadata is missing (\texttt{Unknown}) for 189 of the 200
items. The evaluation therefore probes only a single, narrow slice of the
intended input distribution. A homogeneous test set both reduces the variety of
coaching behaviors exercised and likely contributes to the ceiling effect, since
``partially correct'' answers tend to elicit similarly-structured feedback that
the judge scores uniformly. Rebuilding the test set to span the full range of
answer types and difficulties is the most important step toward a fairer and
more sensitive comparison.

\xhdr{Implications.} Two methodological lessons follow. First, an LLM judge that
saturates near the top of its scale is a poor instrument for measuring
incremental alignment gains; a more discriminative protocol---a finer or
forced-spread rubric, or direct pairwise preference (win rate) judging rather
than absolute scoring---would have more statistical power. Second, preference
optimization that is not explicitly grounded against a reference answer can
introduce fluent but false technical claims; pairing DPO with a
factuality-checking signal is a promising direction. We discuss a small-scale
human evaluation to validate the judge in future work.

%
%  Requires (already in main.tex preamble): booktabs, amssymb,
%  \xhdr command, cleveref. Uses table* (full page width) for the wide
%  statistics table, so it must be in a twocolumn document (it is).
%
%  All numbers below are produced by:
%      python eval/statistical_tests.py   -> eval/statistical_tests.json
%      python eval/failure_analysis.py    -> eval/failure_analysis.json
%      python eval/compare_scores.py      -> eval/comparison_summary.json
%  in the ai-interview-coach repository. Regenerate those if data changes.
% ============================================================================

\section{Experimental Results}
\label{sec:results}

\xhdr{Evaluation protocol.} We evaluate both models on a fixed held-out test
set of 200 examples. The test questions and the simulated student answers are
identical for the baseline and the trained model, so any score difference is
attributable only to the coaching model. For each example, the judge model
(Qwen3.5-9B) scores the feedback on five dimensions (Technical Correctness,
Specificity, Helpfulness, Actionability, Interview Coaching Quality), each on a
1--20 scale; the overall score is their mean. We report per-dimension means and,
because the same 200 items are scored by both models, we treat the two score
vectors as \emph{paired} samples and test the difference with a paired
$t$-test, the Wilcoxon signed-rank test (robust to the non-normal,
ceiling-bound score distribution), and an exact binomial sign test on the
non-tie items. We also report Cohen's $d$ for paired samples as an effect size.

\xhdr{Aggregate results.} As shown in \cref{tab:results}, the SFT+DPO model
attains a marginally higher mean overall score than the baseline
($17.03$ vs.\ $16.86$, $\Delta=+0.17$ on the 1--20 scale). The improvement is
consistent in direction across all five dimensions, but small in magnitude
(per-dimension $\Delta$ between $+0.14$ and $+0.19$).

\xhdr{Statistical significance.} The improvement is \emph{not} statistically
significant. For the overall score, the paired $t$-test gives $t=0.81$,
$p=0.42$, and the Wilcoxon signed-rank test gives $p=0.70$; the effect size is
negligible (Cohen's $d=0.057$). No individual dimension reaches significance
(all $p>0.38$; see \cref{tab:stats}). On an item-by-item basis the trained
model wins on 48 examples, loses on 40, and ties on 112; the sign test over the
88 non-tie items is also non-significant ($p=0.46$). We therefore cannot reject
the null hypothesis that SFT+DPO leaves feedback quality unchanged under this
judge and rubric.

\begin{table}[t]
  \centering
  \small
  \begin{tabular}{@{}lccc@{}}
    \toprule
    Dimension & Baseline & SFT+DPO & $\Delta$ \\
    \midrule
    Technical Correctness        & 16.79 & 16.92 & $+0.14$ \\
    Specificity                  & 16.88 & 17.06 & $+0.19$ \\
    Helpfulness                  & 16.88 & 17.05 & $+0.18$ \\
    Actionability                & 16.89 & 17.05 & $+0.17$ \\
    Interview Coaching Quality   & 16.88 & 17.06 & $+0.19$ \\
    \midrule
    \textbf{Overall Score}       & \textbf{16.86} & \textbf{17.03} & $\mathbf{+0.17}$ \\
    \bottomrule
  \end{tabular}
  \caption{Mean judge scores (1--20) on the 200-example test set. The trained
  model is consistently but marginally higher across all dimensions.}
  \label{tab:results}
\end{table}

\begin{table*}[t]
  \centering
  \small
  \begin{tabular}{@{}lcccccc@{}}
    \toprule
    Dimension & Baseline (SD) & SFT+DPO (SD) & $\Delta$ & Paired $t$ ($p$) & Wilcoxon $p$ & Cohen's $d$ \\
    \midrule
    Technical Correctness        & 16.79 (2.40) & 16.92 (1.92) & $+0.14$ & $0.63\ (0.53)$ & $0.73$ & $0.044$ \\
    Specificity                  & 16.88 (2.37) & 17.06 (1.82) & $+0.19$ & $0.88\ (0.38)$ & $0.65$ & $0.062$ \\
    Helpfulness                  & 16.88 (2.37) & 17.05 (1.85) & $+0.18$ & $0.83\ (0.41)$ & $0.69$ & $0.058$ \\
    Actionability                & 16.89 (2.38) & 17.05 (1.84) & $+0.17$ & $0.78\ (0.44)$ & $0.81$ & $0.055$ \\
    Interview Coaching Quality   & 16.88 (2.37) & 17.06 (1.82) & $+0.19$ & $0.88\ (0.38)$ & $0.69$ & $0.062$ \\
    \midrule
    \textbf{Overall Score}       & \textbf{16.86 (2.37)} & \textbf{17.03 (1.85)} & $\mathbf{+0.17}$ & $\mathbf{0.81\ (0.42)}$ & $\mathbf{0.70}$ & $\mathbf{0.057}$ \\
    \bottomrule
  \end{tabular}
  \caption{Paired statistical comparison of baseline vs.\ SFT+DPO over the 200
  test examples. No dimension reaches significance at $\alpha=0.05$, and all
  effect sizes are negligible ($d<0.07$). Win/loss/tie on the overall score is
  48/40/112; the sign test over non-tie items gives $p=0.46$.}
  \label{tab:stats}
\end{table*}

\section{Discussion}
\label{sec:discussion}

The absence of a significant difference is itself an informative result. We
identify three contributing factors and support them with a per-item failure
analysis.

\xhdr{Ceiling effect.} The judge already rates the baseline very highly,
leaving little headroom for improvement. 65\% of baseline outputs score
$\geq 18/20$, and on 191 of 200 items \emph{both} models score $\geq 15$. When
the starting point is already near the top of the scale, even a genuinely
better model cannot move the mean by much.

\xhdr{Coarse judge resolution.} The judge's scores are effectively discrete,
clustering on a handful of values (predominantly 15 and 18), and the five
rubric dimensions almost always move together (their per-item scores are nearly
identical). This low resolution injects noise and masks the fine-grained
differences that distinguish good coaching feedback from excellent coaching
feedback.

\xhdr{Offsetting gains and regressions.} The flat average hides real movement
in both directions: 48 items improve, 40 regress, and these roughly cancel. The
trained model's largest wins are recoveries of \emph{catastrophic} baseline
failures---three items jump from 1 to 18 (\texttt{test\_0001},
\texttt{test\_0073}, \texttt{test\_0119}) by removing a baseline hallucination
that claimed the student had missed a concept actually present in the answer.
However, the trained model also introduces \emph{new} technical errors that the
baseline did not make: in the most severe regression, \texttt{test\_0194} drops
from 18 to 1 because the feedback falsely asserts that Java's
\texttt{@Override} annotation throws a runtime \texttt{java.lang.Error}.
Categorizing the judge's stated reason on the 40 regressions, 4 are clear
technical inaccuracies and the remainder are smaller, mostly within the judge's
own scoring noise (e.g.\ slightly less complete or less specific feedback).

\xhdr{Test-set diversity limitation.} A confound limits what this experiment can
demonstrate. Although our methodology defines five student-answer types
(correct-but-incomplete, partially correct, incorrect, too vague, verbose but
unfocused), in the released data all 200 test items (and all training items)
carry the label \texttt{partially\_correct}, and the question
category/difficulty metadata is missing (\texttt{Unknown}) for 189 of the 200
items. The evaluation therefore probes only a single, narrow slice of the
intended input distribution. A homogeneous test set both reduces the variety of
coaching behaviors exercised and likely contributes to the ceiling effect, since
``partially correct'' answers tend to elicit similarly-structured feedback that
the judge scores uniformly. Rebuilding the test set to span the full range of
answer types and difficulties is the most important step toward a fairer and
more sensitive comparison.

\xhdr{Implications.} Two methodological lessons follow. First, an LLM judge that
saturates near the top of its scale is a poor instrument for measuring
incremental alignment gains; a more discriminative protocol---a finer or
forced-spread rubric, or direct pairwise preference (win rate) judging rather
than absolute scoring---would have more statistical power. Second, preference
optimization that is not explicitly grounded against a reference answer can
introduce fluent but false technical claims; pairing DPO with a
factuality-checking signal is a promising direction. We discuss a small-scale
human evaluation to validate the judge in future work.

%
%  Requires (already in main.tex preamble): booktabs, amssymb,
%  \xhdr command, cleveref. Uses table* (full page width) for the wide
%  statistics table, so it must be in a twocolumn document (it is).
%
%  All numbers below are produced by:
%      python eval/statistical_tests.py   -> eval/statistical_tests.json
%      python eval/failure_analysis.py    -> eval/failure_analysis.json
%      python eval/compare_scores.py      -> eval/comparison_summary.json
%  in the ai-interview-coach repository. Regenerate those if data changes.
% ============================================================================

\section{Experimental Results}
\label{sec:results}

\xhdr{Evaluation protocol.} We evaluate both models on a fixed held-out test
set of 200 examples. The test questions and the simulated student answers are
identical for the baseline and the trained model, so any score difference is
attributable only to the coaching model. For each example, the judge model
(Qwen3.5-9B) scores the feedback on five dimensions (Technical Correctness,
Specificity, Helpfulness, Actionability, Interview Coaching Quality), each on a
1--20 scale; the overall score is their mean. We report per-dimension means and,
because the same 200 items are scored by both models, we treat the two score
vectors as \emph{paired} samples and test the difference with a paired
$t$-test, the Wilcoxon signed-rank test (robust to the non-normal,
ceiling-bound score distribution), and an exact binomial sign test on the
non-tie items. We also report Cohen's $d$ for paired samples as an effect size.

\xhdr{Aggregate results.} As shown in \cref{tab:results}, the SFT+DPO model
attains a marginally higher mean overall score than the baseline
($17.03$ vs.\ $16.86$, $\Delta=+0.17$ on the 1--20 scale). The improvement is
consistent in direction across all five dimensions, but small in magnitude
(per-dimension $\Delta$ between $+0.14$ and $+0.19$).

\xhdr{Statistical significance.} The improvement is \emph{not} statistically
significant. For the overall score, the paired $t$-test gives $t=0.81$,
$p=0.42$, and the Wilcoxon signed-rank test gives $p=0.70$; the effect size is
negligible (Cohen's $d=0.057$). No individual dimension reaches significance
(all $p>0.38$; see \cref{tab:stats}). On an item-by-item basis the trained
model wins on 48 examples, loses on 40, and ties on 112; the sign test over the
88 non-tie items is also non-significant ($p=0.46$). We therefore cannot reject
the null hypothesis that SFT+DPO leaves feedback quality unchanged under this
judge and rubric.

\begin{table}[t]
  \centering
  \small
  \begin{tabular}{@{}lccc@{}}
    \toprule
    Dimension & Baseline & SFT+DPO & $\Delta$ \\
    \midrule
    Technical Correctness        & 16.79 & 16.92 & $+0.14$ \\
    Specificity                  & 16.88 & 17.06 & $+0.19$ \\
    Helpfulness                  & 16.88 & 17.05 & $+0.18$ \\
    Actionability                & 16.89 & 17.05 & $+0.17$ \\
    Interview Coaching Quality   & 16.88 & 17.06 & $+0.19$ \\
    \midrule
    \textbf{Overall Score}       & \textbf{16.86} & \textbf{17.03} & $\mathbf{+0.17}$ \\
    \bottomrule
  \end{tabular}
  \caption{Mean judge scores (1--20) on the 200-example test set. The trained
  model is consistently but marginally higher across all dimensions.}
  \label{tab:results}
\end{table}

\begin{table*}[t]
  \centering
  \small
  \begin{tabular}{@{}lcccccc@{}}
    \toprule
    Dimension & Baseline (SD) & SFT+DPO (SD) & $\Delta$ & Paired $t$ ($p$) & Wilcoxon $p$ & Cohen's $d$ \\
    \midrule
    Technical Correctness        & 16.79 (2.40) & 16.92 (1.92) & $+0.14$ & $0.63\ (0.53)$ & $0.73$ & $0.044$ \\
    Specificity                  & 16.88 (2.37) & 17.06 (1.82) & $+0.19$ & $0.88\ (0.38)$ & $0.65$ & $0.062$ \\
    Helpfulness                  & 16.88 (2.37) & 17.05 (1.85) & $+0.18$ & $0.83\ (0.41)$ & $0.69$ & $0.058$ \\
    Actionability                & 16.89 (2.38) & 17.05 (1.84) & $+0.17$ & $0.78\ (0.44)$ & $0.81$ & $0.055$ \\
    Interview Coaching Quality   & 16.88 (2.37) & 17.06 (1.82) & $+0.19$ & $0.88\ (0.38)$ & $0.69$ & $0.062$ \\
    \midrule
    \textbf{Overall Score}       & \textbf{16.86 (2.37)} & \textbf{17.03 (1.85)} & $\mathbf{+0.17}$ & $\mathbf{0.81\ (0.42)}$ & $\mathbf{0.70}$ & $\mathbf{0.057}$ \\
    \bottomrule
  \end{tabular}
  \caption{Paired statistical comparison of baseline vs.\ SFT+DPO over the 200
  test examples. No dimension reaches significance at $\alpha=0.05$, and all
  effect sizes are negligible ($d<0.07$). Win/loss/tie on the overall score is
  48/40/112; the sign test over non-tie items gives $p=0.46$.}
  \label{tab:stats}
\end{table*}

\section{Discussion}
\label{sec:discussion}

The absence of a significant difference is itself an informative result. We
identify three contributing factors and support them with a per-item failure
analysis.

\xhdr{Ceiling effect.} The judge already rates the baseline very highly,
leaving little headroom for improvement. 65\% of baseline outputs score
$\geq 18/20$, and on 191 of 200 items \emph{both} models score $\geq 15$. When
the starting point is already near the top of the scale, even a genuinely
better model cannot move the mean by much.

\xhdr{Coarse judge resolution.} The judge's scores are effectively discrete,
clustering on a handful of values (predominantly 15 and 18), and the five
rubric dimensions almost always move together (their per-item scores are nearly
identical). This low resolution injects noise and masks the fine-grained
differences that distinguish good coaching feedback from excellent coaching
feedback.

\xhdr{Offsetting gains and regressions.} The flat average hides real movement
in both directions: 48 items improve, 40 regress, and these roughly cancel. The
trained model's largest wins are recoveries of \emph{catastrophic} baseline
failures---three items jump from 1 to 18 (\texttt{test\_0001},
\texttt{test\_0073}, \texttt{test\_0119}) by removing a baseline hallucination
that claimed the student had missed a concept actually present in the answer.
However, the trained model also introduces \emph{new} technical errors that the
baseline did not make: in the most severe regression, \texttt{test\_0194} drops
from 18 to 1 because the feedback falsely asserts that Java's
\texttt{@Override} annotation throws a runtime \texttt{java.lang.Error}.
Categorizing the judge's stated reason on the 40 regressions, 4 are clear
technical inaccuracies and the remainder are smaller, mostly within the judge's
own scoring noise (e.g.\ slightly less complete or less specific feedback).

\xhdr{Test-set diversity limitation.} A confound limits what this experiment can
demonstrate. Although our methodology defines five student-answer types
(correct-but-incomplete, partially correct, incorrect, too vague, verbose but
unfocused), in the released data all 200 test items (and all training items)
carry the label \texttt{partially\_correct}, and the question
category/difficulty metadata is missing (\texttt{Unknown}) for 189 of the 200
items. The evaluation therefore probes only a single, narrow slice of the
intended input distribution. A homogeneous test set both reduces the variety of
coaching behaviors exercised and likely contributes to the ceiling effect, since
``partially correct'' answers tend to elicit similarly-structured feedback that
the judge scores uniformly. Rebuilding the test set to span the full range of
answer types and difficulties is the most important step toward a fairer and
more sensitive comparison.

\xhdr{Implications.} Two methodological lessons follow. First, an LLM judge that
saturates near the top of its scale is a poor instrument for measuring
incremental alignment gains; a more discriminative protocol---a finer or
forced-spread rubric, or direct pairwise preference (win rate) judging rather
than absolute scoring---would have more statistical power. Second, preference
optimization that is not explicitly grounded against a reference answer can
introduce fluent but false technical claims; pairing DPO with a
factuality-checking signal is a promising direction. We discuss a small-scale
human evaluation to validate the judge in future work.

%
%  Requires (already in main.tex preamble): booktabs, amssymb,
%  \xhdr command, cleveref. Uses table* (full page width) for the wide
%  statistics table, so it must be in a twocolumn document (it is).
%
%  All numbers below are produced by:
%      python eval/statistical_tests.py   -> eval/statistical_tests.json
%      python eval/failure_analysis.py    -> eval/failure_analysis.json
%      python eval/compare_scores.py      -> eval/comparison_summary.json
%  in the ai-interview-coach repository. Regenerate those if data changes.
% ============================================================================

\section{Experimental Results}
\label{sec:results}

\xhdr{Evaluation protocol.} We evaluate both models on a fixed held-out test
set of 200 examples. The test questions and the simulated student answers are
identical for the baseline and the trained model, so any score difference is
attributable only to the coaching model. For each example, the judge model
(Qwen3.5-9B) scores the feedback on five dimensions (Technical Correctness,
Specificity, Helpfulness, Actionability, Interview Coaching Quality), each on a
1--20 scale; the overall score is their mean. We report per-dimension means and,
because the same 200 items are scored by both models, we treat the two score
vectors as \emph{paired} samples and test the difference with a paired
$t$-test, the Wilcoxon signed-rank test (robust to the non-normal,
ceiling-bound score distribution), and an exact binomial sign test on the
non-tie items. We also report Cohen's $d$ for paired samples as an effect size.

\xhdr{Aggregate results.} As shown in \cref{tab:results}, the SFT+DPO model
attains a marginally higher mean overall score than the baseline
($17.03$ vs.\ $16.86$, $\Delta=+0.17$ on the 1--20 scale). The improvement is
consistent in direction across all five dimensions, but small in magnitude
(per-dimension $\Delta$ between $+0.14$ and $+0.19$).

\xhdr{Statistical significance.} The improvement is \emph{not} statistically
significant. For the overall score, the paired $t$-test gives $t=0.81$,
$p=0.42$, and the Wilcoxon signed-rank test gives $p=0.70$; the effect size is
negligible (Cohen's $d=0.057$). No individual dimension reaches significance
(all $p>0.38$; see \cref{tab:stats}). On an item-by-item basis the trained
model wins on 48 examples, loses on 40, and ties on 112; the sign test over the
88 non-tie items is also non-significant ($p=0.46$). We therefore cannot reject
the null hypothesis that SFT+DPO leaves feedback quality unchanged under this
judge and rubric.

\begin{table}[t]
  \centering
  \small
  \begin{tabular}{@{}lccc@{}}
    \toprule
    Dimension & Baseline & SFT+DPO & $\Delta$ \\
    \midrule
    Technical Correctness        & 16.79 & 16.92 & $+0.14$ \\
    Specificity                  & 16.88 & 17.06 & $+0.19$ \\
    Helpfulness                  & 16.88 & 17.05 & $+0.18$ \\
    Actionability                & 16.89 & 17.05 & $+0.17$ \\
    Interview Coaching Quality   & 16.88 & 17.06 & $+0.19$ \\
    \midrule
    \textbf{Overall Score}       & \textbf{16.86} & \textbf{17.03} & $\mathbf{+0.17}$ \\
    \bottomrule
  \end{tabular}
  \caption{Mean judge scores (1--20) on the 200-example test set. The trained
  model is consistently but marginally higher across all dimensions.}
  \label{tab:results}
\end{table}

\begin{table*}[t]
  \centering
  \small
  \begin{tabular}{@{}lcccccc@{}}
    \toprule
    Dimension & Baseline (SD) & SFT+DPO (SD) & $\Delta$ & Paired $t$ ($p$) & Wilcoxon $p$ & Cohen's $d$ \\
    \midrule
    Technical Correctness        & 16.79 (2.40) & 16.92 (1.92) & $+0.14$ & $0.63\ (0.53)$ & $0.73$ & $0.044$ \\
    Specificity                  & 16.88 (2.37) & 17.06 (1.82) & $+0.19$ & $0.88\ (0.38)$ & $0.65$ & $0.062$ \\
    Helpfulness                  & 16.88 (2.37) & 17.05 (1.85) & $+0.18$ & $0.83\ (0.41)$ & $0.69$ & $0.058$ \\
    Actionability                & 16.89 (2.38) & 17.05 (1.84) & $+0.17$ & $0.78\ (0.44)$ & $0.81$ & $0.055$ \\
    Interview Coaching Quality   & 16.88 (2.37) & 17.06 (1.82) & $+0.19$ & $0.88\ (0.38)$ & $0.69$ & $0.062$ \\
    \midrule
    \textbf{Overall Score}       & \textbf{16.86 (2.37)} & \textbf{17.03 (1.85)} & $\mathbf{+0.17}$ & $\mathbf{0.81\ (0.42)}$ & $\mathbf{0.70}$ & $\mathbf{0.057}$ \\
    \bottomrule
  \end{tabular}
  \caption{Paired statistical comparison of baseline vs.\ SFT+DPO over the 200
  test examples. No dimension reaches significance at $\alpha=0.05$, and all
  effect sizes are negligible ($d<0.07$). Win/loss/tie on the overall score is
  48/40/112; the sign test over non-tie items gives $p=0.46$.}
  \label{tab:stats}
\end{table*}

\section{Discussion}
\label{sec:discussion}

The absence of a significant difference is itself an informative result. We
identify three contributing factors and support them with a per-item failure
analysis.

\xhdr{Ceiling effect.} The judge already rates the baseline very highly,
leaving little headroom for improvement. 65\% of baseline outputs score
$\geq 18/20$, and on 191 of 200 items \emph{both} models score $\geq 15$. When
the starting point is already near the top of the scale, even a genuinely
better model cannot move the mean by much.

\xhdr{Coarse judge resolution.} The judge's scores are effectively discrete,
clustering on a handful of values (predominantly 15 and 18), and the five
rubric dimensions almost always move together (their per-item scores are nearly
identical). This low resolution injects noise and masks the fine-grained
differences that distinguish good coaching feedback from excellent coaching
feedback.

\xhdr{Offsetting gains and regressions.} The flat average hides real movement
in both directions: 48 items improve, 40 regress, and these roughly cancel. The
trained model's largest wins are recoveries of \emph{catastrophic} baseline
failures---three items jump from 1 to 18 (\texttt{test\_0001},
\texttt{test\_0073}, \texttt{test\_0119}) by removing a baseline hallucination
that claimed the student had missed a concept actually present in the answer.
However, the trained model also introduces \emph{new} technical errors that the
baseline did not make: in the most severe regression, \texttt{test\_0194} drops
from 18 to 1 because the feedback falsely asserts that Java's
\texttt{@Override} annotation throws a runtime \texttt{java.lang.Error}.
Categorizing the judge's stated reason on the 40 regressions, 4 are clear
technical inaccuracies and the remainder are smaller, mostly within the judge's
own scoring noise (e.g.\ slightly less complete or less specific feedback).

\xhdr{Test-set diversity limitation.} A confound limits what this experiment can
demonstrate. Although our methodology defines five student-answer types
(correct-but-incomplete, partially correct, incorrect, too vague, verbose but
unfocused), in the released data all 200 test items (and all training items)
carry the label \texttt{partially\_correct}, and the question
category/difficulty metadata is missing (\texttt{Unknown}) for 189 of the 200
items. The evaluation therefore probes only a single, narrow slice of the
intended input distribution. A homogeneous test set both reduces the variety of
coaching behaviors exercised and likely contributes to the ceiling effect, since
``partially correct'' answers tend to elicit similarly-structured feedback that
the judge scores uniformly. Rebuilding the test set to span the full range of
answer types and difficulties is the most important step toward a fairer and
more sensitive comparison.

\xhdr{Implications.} Two methodological lessons follow. First, an LLM judge that
saturates near the top of its scale is a poor instrument for measuring
incremental alignment gains; a more discriminative protocol---a finer or
forced-spread rubric, or direct pairwise preference (win rate) judging rather
than absolute scoring---would have more statistical power. Second, preference
optimization that is not explicitly grounded against a reference answer can
introduce fluent but false technical claims; pairing DPO with a
factuality-checking signal is a promising direction. We discuss a small-scale
human evaluation to validate the judge in future work.
